\cleardoublepage
\sectionnonum{作者简历}

\begingroup
\footnotesize
\setlength{\parindent}{0pt}
\setlength{\parskip}{0pt}
\setlength{\tabcolsep}{2.5pt}
\renewcommand{\arraystretch}{1.0}

\textbf{基本信息}

\begin{tabularx}{\textwidth}{@{}lXlX@{}}
    姓名：于瑞梓 & & 年龄：22 岁 & \\
    专业：计算机科学与技术 & & 学历：本科 & \\
    手机：17348548980 & & 邮箱：yuruizi126@zju.edu.cn & \\
\end{tabularx}

\vspace{0.1em}
\textbf{教育经历}

\begin{tabularx}{\textwidth}{@{}lX@{}}
    2022.09--至今 & 浙江大学，计算机科学与技术专业，本科在读。
\end{tabularx}

\vspace{0.1em}
\textbf{科研与项目经历}

\begin{tabularx}{\textwidth}{@{}p{0.17\textwidth}X@{}}
    2025.04--2025.05 &
    \textbf{SportsEye: Human-Inspired Active Perception for Fine-Grained Sports Video Understanding}。参与构建面向体育视频细粒度理解的时空推理数据集 SportsTrackQA-1M 与评测基准 SportsTrackBench，负责部分数据构建流程代码，实现视频切片、人物追踪与掩码生成、关键帧采样、辅助信息标注和 prompt 构建等标准化处理；基于 MOT Tracking、GroundingDINO 与 SAM2 生成跨帧轨迹和实例级 mask 标注，并参与体育视频时空推理问答对生成流程设计。 \\
    2024.06--2024.09 &
    \textbf{基于深度学习的考古报告版面分析与信息提取工具}（良渚博物院合作项目）。面向考古文献图像整理需求，设计并实现基于深度学习的文献图像解析系统，支持 PDF 报告中器物图像的自动检测、分割、信息提取与交互式修正；使用 YOLOv10 完成器物目标检测，结合 Segment Anything Model 实现交互式精细分割，并使用 PaddleOCR 识别图注与编号，完成图像自动命名和结构化整理。项目产出国家发明专利 1 项，并获 2024 中国国际大学生创新大赛浙江省金奖、全国铜奖。 \\
\end{tabularx}

\vspace{0.1em}
\textbf{获奖情况}

\begin{tabularx}{\textwidth}{@{}p{0.17\textwidth}X@{}}
    2024.07--2024.08 & 烛龙科技：中国国际大学生创新大赛（2024）全国金奖。 \\
    2024.01--2024.08 & 智载千古：中国国际大学生创新大赛（2024）全国铜奖。 \\
    2024.01--2024.04 & 浙江大学第十六届“蒲公英”大学生创业大赛一等奖。 \\
    2024.04--2024.07 & 第九届全国大学生生物医学工程创新设计竞赛三等奖。 \\
    2024.08--2024.10 & 浙江大学第三届“智云杯”学生智慧校园应用创新大赛三等奖。 \\
\end{tabularx}

\vspace{0.1em}
\textbf{学生工作与综合能力}

\begin{tabularx}{\textwidth}{@{}p{0.17\textwidth}X@{}}
    2024.05--2025.05 & 浙江大学启真交叉学科创新创业实验室软件团队管理层，参与并领导团队项目开发。 \\
    2023.07--2024.06 & 浙江大学 D.F.M. 学生街舞社 locking 舞队队长，组织日常训练与校际演出。 \\
    2022.09--2023.06 & 浙江大学云峰融媒体中心干事，运营校园公众号并参与校园新闻摄影。 \\
    其他能力 & 通过 CET-6、CET-4，具有红十字会急救员证，获浙江大学文体活动标兵称号。 \\
\end{tabularx}

\endgroup
